\section{The actual clipped Chebyshev filter}

\begin{lemma}[Spatial decomposition and norm of the clipped Hamiltonian]
    \label{lem:aklv_clipped_hamiltonian_data}
    \lean{TNLean.MPS.AreaLaw.aklv_clipped_hamiltonian_decomposition,
        TNLean.MPS.AreaLaw.aklv_clipped_hamiltonian_posSemidef_norm_le}
    \leanok
    \uses{def:aklv_normalization_regions,def:aklv_two_block_truncation}
    Let $s\ge6$, $a+s+1\le L$, and let $0\le h_i\le\Id$ act on
    $\{i,i+1\}$. Let $e_0,e_1,e_2$ be lower bounds for the three reference
    Hamiltonians, with $e_1\ge0$, and let $t\ge0$.
    The actual two-exterior truncation $B$ satisfies
    \begin{align*}
        0\le B,\qquad \|B\|\le2t+s.
    \end{align*}
    It can be written as
    \begin{align*}
        B=H'_L+\sum_{j=0}^{s-3}h_{a+1+j}+H'_R,
    \end{align*}
    where $H'_L$ acts on $[0,a+2)$ and $H'_R$ acts on $[a+s-1,L)$. % chktex 44
    This window has width $s-2$ and the same central cut, since
    \begin{align*}
        a+1+\left\lfloor\frac{s-2}{2}\right\rfloor
        =a+\left\lfloor\frac{s}{2}\right\rfloor.
    \end{align*}
\end{lemma}
\begin{proof}
    \leanok
    The two clipped exterior sums have ceilings $e_0+t$ and $e_2+t$.
    The $s-2$ retained bonds have sum at most $(s-2)\Id$.
    After subtracting $e_0+e_1+e_2$, their sum is at most
    $(2t+s-2-e_1)\Id\le(2t+s)\Id$.
    Each centered clipped region is positive because its clipping ceiling
    is at least its regional lower bound; the connecting bonds are positive.
    The norm estimate follows. Functional calculus preserves the exterior
    supports. Placing the scalar shift in the left exterior gives the displayed
    decomposition of the same operator.
\end{proof}

\begin{lemma}[A clipped physical filter]
    \label{thm:aklv_clipped_physical_filter}
    \lean{TNLean.MPS.AreaLaw.exists_aklv_clipped_chebyshev_filter}
    \leanok
    \uses{thm:aklv_physical_robustness,lem:aklv_clipped_hamiltonian_data}
    Under the physical hypotheses of
    Lemma~\ref{thm:aklv_physical_robustness}, choose the three actual
    reference minima and the normalized original ground energy $E_0$.
    For every
    \begin{align*}
        t\ge\frac{1024}{g}\left(\frac{E_0}{g}+1\right),
    \end{align*}
    the actual two-exterior truncation $B$ has a unit ground vector $\phi$
    with energy $E\ge0$, gap at least $9g/10$, and
    \begin{align*}
        \|\phi-\Omega\|^2\le e^{-t/16}/100.
    \end{align*}
    For each integer $\ell\ge0$ there is a real polynomial $p$ of degree
    $\ell$, with $p(E)=1$, such that $K=p(B)$ fixes $\phi$ and
    \begin{align*}
        \|K-|\phi\rangle\langle\phi|\|
         & \le A,
         & A                           & =2\exp\left(-2\ell\sqrt{\frac{9g}{10(2t+s)}}\right), \\
        \operatorname{OSR}_{a+\lfloor s/2\rfloor}(K)
         & \le\bigl(q(\ell+1)\bigr)^ {
                32(\lfloor\ell/s\rfloor+\lfloor\sqrt\ell\rfloor+1)}.
    \end{align*}
    Moreover, $K$ preserves $\phi^\perp$ and
    $\|K\psi\|^2\le A^2$ for every unit $\psi\perp\phi$.
    The vector $\phi$ and its energy are chosen before $\ell$.
    This constructs the physical filter used after
    \cite[Theorem~6.1]{Arad2013AreaLaw}, for the stated window range;
    it does not assert the later varying-truncation entropy construction.
\end{lemma}
\begin{proof}
    \leanok
    \uses{thm:aklv_physical_chebyshev_filter}
    The actual canonical eigenvalues obey a gap of $9g/10$ above the
    distinguished ground index. Their complete orthonormal basis therefore
    gives
    \begin{align*}
        B-E\Id\ge\frac{9g}{10}
        (\Id-|\phi\rangle\langle\phi|).
    \end{align*}
    The preceding lemma supplies the ceiling $2t+s$ and the shorter supported
    window at the original central cut. Apply the Chebyshev construction and
    the sharp polynomial rank estimate to this same matrix. Its exact
    denominator is $2t+s-E$; replacing it by $2t+s$ weakens the error bound
    because $E\ge0$. For $s\ge6$, division with remainder gives
    $\lfloor\ell/(s-2)\rfloor\le2\lfloor\ell/s\rfloor+1$, which
    gives the displayed exponent. The squared shrinking estimate follows on
    $\phi^\perp$.
\end{proof}

\begin{lemma}[The clipped filter above a logarithmic threshold]
    \label{cor:aklv_logarithmic_clipped_filter}
    \lean{TNLean.MPS.AreaLaw.exists_aklv_logarithmic_clipped_chebyshev_filter}
    \leanok
    \uses{def:aklv_logarithmic_threshold,thm:aklv_clipped_physical_filter}
    Under the physical hypotheses of
    Lemma~\ref{thm:aklv_clipped_physical_filter}, one has $q\ge2$ and
    $g\le1$. The regional minima can be chosen before $t$ so that every
    conclusion of that theorem holds already for
    \begin{align*}
        t\ge T(g)=1024+16\log(1/g).
    \end{align*}
    In particular, the actual ground vector, projector gap, spectral ceiling,
    and same-cut polynomial rank bound require no supplied auxiliary data.
\end{lemma}
\begin{proof}
    \leanok
    \uses{thm:aklv_logarithmic_robustness,
        lem:aklv_clipped_hamiltonian_data,thm:aklv_physical_chebyshev_filter}
    Use the stronger logarithmic robustness theorem to construct the
    actual ground vector. The norm, spatial decomposition and Chebyshev
    argument of Lemma~\ref{thm:aklv_clipped_physical_filter} then apply
    unchanged to the same clipped matrix. The polynomial-threshold
    version follows from $T(g)\le1024/g$ and $E_0\ge0$.
\end{proof}
